\documentclass[11pt,a4paper]{article}
\usepackage[utf8]{inputenc}
\usepackage[T1]{fontenc}
\usepackage[portuguese]{babel}
\usepackage{xcolor}
\usepackage{hyperref}
\usepackage{geometry}
\usepackage{tcolorbox}

\geometry{margin=2.5cm}

\definecolor{primary}{RGB}{41, 128, 185}
\definecolor{secondary}{RGB}{44, 62, 80}
\definecolor{codebg}{RGB}{240, 240, 240}

\hypersetup{colorlinks=true, linkcolor=primary, urlcolor=primary}

\title{\color{secondary}\Huge\textbf{Exercício 5.3: Log app, a Edição Service Mesh}}
\author{\color{primary}DevOps com Kubernetes -- 2026}
\date{}

\begin{document}
\maketitle

\section*{\color{primary}Instruções do Exercício}
\begin{tcolorbox}[colback=blue!5!white,colframe=blue!75!black,title=Exercício 5.3: Log app, a Edição Service Mesh]
Implemente a aplicação Log output na malha de serviço (\textit{service mesh}). 

- Estenda a aplicação desenvolvendo um novo serviço chamado \textit{greeter}, capaz de responder a uma requisição HTTP GET com uma saudação ('Hello').

- A aplicação \textit{log output} deve passar a consumir o serviço \textit{greeter} via rede e exibir a saudação na tela juntamente com as saídas dos logs regulares.

- Implante duas versões diferentes do \textit{greeter} (v1 e v2) e configure as regras de roteamento do Istio (como o \textit{HTTPRoute}) para dividir o tráfego: uma versão deve receber 75\% das requisições e a outra os 25\% restantes.
\end{tcolorbox}

\end{document}
